% ------------------------------------------------------------
% Section 4 -- Experimental Setup
% ------------------------------------------------------------
\section{Experimental Setup}
\label{sec:setup}

\subsection{Data}
\label{sec:data}

Mayo Clinic Low Dose CT dataset (AAPM Low Dose CT Grand Challenge;~\cite{mccollough2017}), accessed and split via the ldct-benchmark tooling~\cite{eulig2024} to guarantee protocol comparability (100p split). Held-out test set of 10 patients, balanced by anatomy: chest C121, C135, C170, C249, C280; abdomen L006, L107, L220, L221, L241. Pilot screening used a disjoint 8/4 train/validation patient subset, with all screening physics metrics computed on validation patients only --- no test leakage. Preprocessing follows the benchmark preset: HU window $-1024$ to $1900$, standardization mean $481.4542$ / std $502.1851$ (offset 1024), $128 \times 128$ training patches, batch size 73. NPS estimation uses central $320 \times 320$ crops. The real-noise reference is the residual $\mathrm{LDCT} - \mathrm{NDCT}$.

\subsection{Matched-budget protocol}
\label{sec:protocol}

Every configuration is trained with an identical budget of 30{,}000 iterations (validation every 1{,}000; learning rate $9.5834 \times 10^{-5}$, constant; model selection by benchmark SSIM on validation). The three decisive arms --- C0, C1h-w0.01, and S4b --- are each trained with three random seeds (0, 1, 2); all other ablation arms use seed 0. A cross-trunk replication additionally trains the same three arms on a ResNet trunk~\cite{he2016} from the same benchmark suite (seed 0), with identical protocol, budget, and loss weights and no re-tuning. We train for 30k iterations because validation metrics plateau well before that point and the residual improvement beyond it is small relative to its compute cost; identical budgets also keep every comparison attribution-clean --- an under-trained model removes less noise and spuriously ``wins'' NPS metrics, so cross-budget comparisons are disallowed by protocol. Loss weights were calibrated in a screening phase (8/4 patients, 8k iterations); only rankings, never values, transfer from screening.

\subsection{Configurations}
\label{sec:configs}

Table~\ref{tab:configs} lists all configurations; all are trained under the identical protocol above.

\begin{table}[t]
\centering
\caption{Configurations.}
\label{tab:configs}
\small
\setlength{\tabcolsep}{3pt}
\begin{tabular}{lllll}
\toprule
ID & Head & $\wnps$ & $\whu$ & Role \\
\midrule
C0 & --- & 0 & 0 & Baseline (MSE only) \\
C2 & --- & 0 & 0.2 & HU-bin loss only \\
C1h & --- & 0.005 & 0.2 & Best loss-only recipe \\
C1h-w0.01 & --- & 0.01 & 0.2 & \textbf{Matched-weight control} \\
S1 & static & 0 & 0 & Head w/o incentive (neutrality) \\
S3b & static, frz-DC 2 & 0.005 & 0.2 & Static head \\
S4 & adaptive, frz-DC 2 & 0.005 & 0.2 & Diagnostic (saturation failure) \\
\textbf{S4b} & adapt.\ + cent.\ 0.05, frz-DC 2 & 0.01 & 0.2 & \textbf{Proposed} \\
\bottomrule
\end{tabular}
\end{table}

\subsection{Metrics}
\label{sec:metrics}

Image quality: PSNR, SSIM, VIF, RMSE\textsubscript{HU}. Physical fidelity (removed noise $= x - \hat{x}$, real noise $= \mathrm{LDCT} - \mathrm{NDCT}$, central $320 \times 320$ crops):

\begin{enumerate}
  \item \textbf{\npsl{} ($\downarrow$)} --- $L_1$ distance between radially averaged log noise-power spectra of removed vs.\ real noise; the primary spectral-fidelity metric.
  \item \textbf{\npsc{} ($\uparrow$)} --- Pearson correlation of the radial NPS curves; spectral shape fidelity independent of magnitude.
  \item \textbf{\stdr{} ($\rightarrow 1$)} --- std of removed noise / std of real noise; below 1 means residual noise left in the image, above 1 means anatomy erased as ``noise''.
  \item \textbf{Per-tissue HU bias ($\rightarrow 0$\hu{})} --- mean signed error within the fixed HU bins of Sec.~\ref{sec:objective}; directly clinically interpretable.
\end{enumerate}

\paragraph{Reporting rule --- always per anatomy.}
Overall aggregates can cancel: C2's overall bone bias of $-2.0$\hu{} hides chest $-18.3$ / abdomen $+14.3$\hu{}. All tables report chest and abdomen separately; overall values are shown for convenience only. Statistical testing: paired per-patient Wilcoxon signed-rank tests ($n = 10$, two-sided; a 10/10 sweep yields the exact minimum $p = 0.002$).
